% ------------------------------------------------------------------
% Title page, how to use, session plan, notation.
% \iffull is true for the Chapters 0-5 notes, false for Chapters 3-5.
% ------------------------------------------------------------------
\begin{titlepage}
\centering
\vspace*{3cm}
{\scshape\large Study notes\par}
\vspace{0.6cm}
{\Huge\bfseries Metric Spaces \& Topology\par}
\vspace{0.5cm}
\iffull
{\Large Chapters 0--5\par}
\else
{\Large Chapters 3--5: Continuity, Completeness, Compactness\par}
\fi
\vspace{1.2cm}
\rule{0.6\textwidth}{0.4pt}\par
\vspace{0.6cm}
{\large Tuesday 10 November \quad 2:00--7:00\,pm session\par}
\vfill
{\small Definitions, theorems with proofs, worked examples, common mistakes,\\
graded exercises with full solutions, proof templates, a counterexample bank\\
and a glossary of every term needed for the exam.\par}
\vspace{1cm}
\end{titlepage}

\pagenumbering{roman}
\chapter*{How to use these notes}
\addcontentsline{toc}{chapter}{How to use these notes}
\markboth{How to use these notes}{}

Each chapter starts with its \textbf{study slot} from the session plan below and a short list of what to focus on in that slot. Inside each chapter:

\begin{itemize}
  \item \textbf{Intuition} paragraphs give the idea in plain English before the formal statement.
  \item \colorbox{defbg}{\textbf{Definitions}} are worded the way you should state them in an exam.
  \item \colorbox{thmbg}{\textbf{Theorems, lemmas, propositions}} come with proofs of the results most often examined.
  \item \colorbox{exbg}{\textbf{Examples}} and \colorbox{warnbg}{\textbf{Warnings}} cover standard cases and common mistakes.
  \item \textbf{Exercises} are tagged \easy, \medium\ or \hard\ and each has a full solution. Cover the solution and try the exercise first.
\end{itemize}

The appendices contain a \textbf{proof-template toolkit}, a \textbf{counterexample bank}, a \textbf{glossary} and the \textbf{recap checklist} for the 5:35 slot.
\iffull\else
Since these notes start at Chapter~3, Appendix~\ref{app:prereq} summarises everything from Chapters~0--2 that the later chapters use. Read it first if those chapters aren't fresh.
\fi

\medskip
{\small\itshape Chapter numbering follows the usual order for a first course in metric spaces. If your lecture notes order the topics differently, go by the section headings.}

\section*{Session plan: Tuesday 10 November}
\begin{center}
\renewcommand{\arraystretch}{1.15}
\begin{tabular}{@{}lll@{}}
\toprule
\textbf{Time} & \textbf{Task} & \textbf{Where} \\
\midrule
2:00--2:10 & Warm-up: skim the chapter headings & Contents \\
2:10--2:30 & Chapter 0: preliminaries & \iffull Chapter 0\else Appendix~\ref{app:prereq}\fi \\
2:30--3:05 & Chapter 1: metric spaces & \iffull Chapter 1\else Appendix~\ref{app:prereq}\fi \\
3:05--3:15 & Break & \\
3:15--3:50 & Chapter 2: open and closed sets, topologies & \iffull Chapter 2\else Appendix~\ref{app:prereq}\fi \\
3:50--4:25 & Chapter 3: continuity & Chapter 3 \\
4:25--4:40 & Longer break & \\
4:40--5:10 & Chapter 4: completeness & Chapter 4 \\
5:10--5:35 & Chapter 5: compactness & Chapter 5 \\
5:35--5:45 & Recap from memory & Recap checklist \\
5:45--6:00 & Break & \\
6:00--6:50 & Linear systems & Separate notes \\
6:50--7:00 & Wrap-up: what's done, where to start next time & \\
\bottomrule
\end{tabular}
\end{center}

\section*{Notation}
\begin{tabularx}{\textwidth}{@{}lY@{}}
$(X,d)$ & a metric space; $x,y,z\in X$ are points \\
$\Bo(x,r)=\set{y\in X}{d(x,y)<r}$ & open ball of radius $r>0$ \\
$\Bc(x,r)=\set{y\in X}{d(x,y)\le r}$ & closed ball \\
$\inte A,\ \cl A = \overline{A},\ \partial A$ & interior, closure, boundary \\
$A^c = X\setminus A$ & complement \\
$A'$ & derived set (the limit points of $A$) \\
$\seq{x}$, $x_n\to x$ & a sequence; convergence \\
$\N=\{1,2,3,\dots\}$, $\Z,\Q,\R,\C$ & the usual number systems \\
$C[a,b]$ & the continuous functions $[a,b]\to\R$ \\
\end{tabularx}

\tableofcontents
\clearpage
\pagenumbering{arabic}
